\section{Learning Objectives}

After completing this chapter, readers will be able to:
\begin{itemize}
    \item Design and implement automated retraining pipeline architectures
    \item Configure intelligent trigger mechanisms for model retraining
    \item Implement incremental learning and online updating strategies
    \item Handle concept drift and data drift in production systems
    \item Balance model freshness against stability requirements
    \item Build continuous learning systems for real-time adaptation
    \item Automate model maintenance workflows end-to-end
    \item Monitor and optimize retraining costs and performance
\end{itemize}

\section{Introduction}

Machine learning models degrade over time as the world changes. Customer behavior shifts, market conditions evolve, and data distributions drift from training assumptions. Without systematic maintenance, production models become stale, accuracy drops, and business value erodes.

Model maintenance encompasses the complete lifecycle of keeping models relevant and performant:

\begin{itemize}
    \item \textbf{Performance monitoring:} Continuous tracking of model metrics
    \item \textbf{Drift detection:} Identifying when retraining is needed
    \item \textbf{Automated retraining:} Triggering and executing model updates
    \item \textbf{Validation:} Ensuring new models improve upon current ones
    \item \textbf{Deployment:} Safely rolling out updated models
    \item \textbf{Rollback:} Reverting problematic updates quickly
\end{itemize}

This chapter provides production-ready architectures and code for building robust model maintenance systems that keep ML applications healthy without manual intervention.

\section{Automated Retraining Pipeline Architecture}

\subsection{Core Components}

A production retraining pipeline consists of multiple coordinated components working together to maintain model quality.

\begin{lstlisting}[language=Python, caption=Complete Automated Retraining Pipeline]
from typing import Dict, List, Optional, Any, Callable
from dataclasses import dataclass, field
from datetime import datetime, timedelta
from enum import Enum
import pandas as pd
import numpy as np
from pathlib import Path
import logging
import json
import pickle

class RetrainingTrigger(Enum):
    """Triggers that can initiate retraining"""
    PERFORMANCE_DEGRADATION = "performance_degradation"
    DATA_DRIFT = "data_drift"
    CONCEPT_DRIFT = "concept_drift"
    SCHEDULED = "scheduled"
    MANUAL = "manual"
    DATA_VOLUME = "data_volume"
    TIME_BASED = "time_based"

class ModelStatus(Enum):
    """Model lifecycle status"""
    TRAINING = "training"
    VALIDATING = "validating"
    STAGING = "staging"
    PRODUCTION = "production"
    ARCHIVED = "archived"
    FAILED = "failed"

@dataclass
class RetrainingConfig:
    """Configuration for retraining pipeline"""
    model_name: str
    model_type: str

    # Trigger thresholds
    performance_threshold: float = 0.05  # 5% degradation
    drift_threshold: float = 0.15
    min_data_points: int = 1000
    max_days_without_retrain: int = 30

    # Training settings
    validation_split: float = 0.2
    test_split: float = 0.1
    max_training_time_hours: int = 24

    # Deployment settings
    canary_percentage: float = 0.1
    canary_duration_hours: int = 24
    auto_promote: bool = False

    # Resource limits
    max_memory_gb: int = 32
    max_cpu_cores: int = 16
    use_gpu: bool = False

@dataclass
class RetrainingEvent:
    """Record of a retraining event"""
    event_id: str
    timestamp: datetime
    trigger: RetrainingTrigger
    model_version: str
    metrics: Dict[str, float]
    data_stats: Dict[str, Any]
    training_duration: Optional[float] = None
    success: bool = False
    error_message: Optional[str] = None

class DataCollector:
    """
    Collect and manage training data for retraining

    Handles incremental data collection, quality checks,
    and dataset versioning
    """

    def __init__(self, data_path: Path, max_samples: int = 1000000):
        self.data_path = data_path
        self.max_samples = max_samples
        self.logger = logging.getLogger(__name__)

    def collect_new_data(self,
                        since: datetime,
                        quality_checks: bool = True) -> pd.DataFrame:
        """Collect new training data since timestamp"""

        # In production, this would query your data warehouse
        # For demo, simulate data collection
        self.logger.info(f"Collecting data since {since}")

        # Simulate fetching new data
        new_data = self._fetch_from_source(since)

        if quality_checks:
            new_data = self._apply_quality_checks(new_data)

        self.logger.info(f"Collected {len(new_data)} new samples")
        return new_data

    def _fetch_from_source(self, since: datetime) -> pd.DataFrame:
        """
        Fetch data from source.

        Reference implementation: load a local Parquet/CSV snapshot if available;
        otherwise generate a synthetic dataset so the example remains runnable.
        """
        try:
            if self.data_path.exists():
                self.logger.info(f"Loading data from {self.data_path}")
                if self.data_path.suffix == ".parquet":
                    df = pd.read_parquet(self.data_path)
                else:
                    df = pd.read_csv(self.data_path)
            else:
                self.logger.warning(
                    f"{self.data_path} not found. Generating synthetic dataset for demo."
                )
                # Synthetic classification data with drift on timestamp
                n_rows = min(self.max_samples, 10_000)
                rng = np.random.default_rng(42)
                df = pd.DataFrame(
                    rng.normal(size=(n_rows, 10)),
                    columns=[f"feature_{i}" for i in range(10)]
                )
                df["timestamp"] = pd.date_range(end=datetime.now(), periods=n_rows, freq="min")
                df["target"] = (df.sum(axis=1) > 0).astype(int)

            # Filter by timestamp if present
            if "timestamp" in df.columns:
                df["timestamp"] = pd.to_datetime(df["timestamp"], errors="coerce")
                df = df[df["timestamp"] >= since]

            if df.empty:
                raise ValueError("No data available after filtering by timestamp")

            return df

        except Exception as e:
            self.logger.error(f"Failed to fetch data: {e}")
            raise

    def _apply_quality_checks(self, data: pd.DataFrame) -> pd.DataFrame:
        """Apply data quality checks"""

        # Remove duplicates
        initial_count = len(data)
        data = data.drop_duplicates()

        # Remove nulls in critical columns
        data = data.dropna(subset=['target'])

        # Remove outliers (example: using IQR)
        for col in data.select_dtypes(include=[np.number]).columns:
            Q1 = data[col].quantile(0.25)
            Q3 = data[col].quantile(0.75)
            IQR = Q3 - Q1
            data = data[
                (data[col] >= Q1 - 3*IQR) &
                (data[col] <= Q3 + 3*IQR)
            ]

        removed = initial_count - len(data)
        if removed > 0:
            self.logger.warning(f"Removed {removed} samples in quality checks")

        return data

    def get_training_dataset(self,
                           days_back: int = 90,
                           include_validation: bool = True) -> Dict[str, pd.DataFrame]:
        """Get complete training dataset"""

        cutoff_date = datetime.now() - timedelta(days=days_back)
        data = self.collect_new_data(cutoff_date)

        # Limit to max samples
        if len(data) > self.max_samples:
            data = data.sample(n=self.max_samples, random_state=42)

        if not include_validation:
            return {'train': data}

        # Split into train/val
        from sklearn.model_selection import train_test_split
        train, val = train_test_split(data, test_size=0.2, random_state=42)

        return {
            'train': train,
            'validation': val
        }

class PerformanceEvaluator:
    """
    Evaluate model performance and compare versions

    Determines if new model is better than current production model
    """

    def __init__(self, metrics: List[str] = None):
        self.metrics = metrics or ['accuracy', 'precision', 'recall', 'f1']
        self.logger = logging.getLogger(__name__)

    def evaluate_model(self,
                      model: Any,
                      test_data: pd.DataFrame,
                      target_col: str = 'target') -> Dict[str, float]:
        """Evaluate model on test data"""

        X_test = test_data.drop(columns=[target_col])
        y_test = test_data[target_col]

        predictions = model.predict(X_test)

        from sklearn.metrics import (
            accuracy_score, precision_score,
            recall_score, f1_score
        )

        results = {
            'accuracy': accuracy_score(y_test, predictions),
            'precision': precision_score(y_test, predictions, average='weighted'),
            'recall': recall_score(y_test, predictions, average='weighted'),
            'f1': f1_score(y_test, predictions, average='weighted'),
            'samples': len(test_data)
        }

        return results

    def compare_models(self,
                      new_metrics: Dict[str, float],
                      current_metrics: Dict[str, float],
                      primary_metric: str = 'f1',
                      min_improvement: float = 0.01) -> Dict[str, Any]:
        """Compare new model against current production model"""

        comparison = {
            'new_better': False,
            'improvement': {},
            'primary_metric': primary_metric,
            'recommendation': 'keep_current'
        }

        for metric in self.metrics:
            if metric in new_metrics and metric in current_metrics:
                improvement = new_metrics[metric] - current_metrics[metric]
                comparison['improvement'][metric] = improvement

        # Check primary metric
        if primary_metric in comparison['improvement']:
            improvement = comparison['improvement'][primary_metric]

            if improvement > min_improvement:
                comparison['new_better'] = True
                comparison['recommendation'] = 'promote_new'
                self.logger.info(
                    f"New model shows {improvement:.4f} improvement "
                    f"in {primary_metric}"
                )
            elif improvement < -min_improvement:
                comparison['recommendation'] = 'keep_current'
                self.logger.warning(
                    f"New model performs worse by {abs(improvement):.4f} "
                    f"in {primary_metric}"
                )
            else:
                comparison['recommendation'] = 'neutral'

        return comparison

class AutomatedRetrainingPipeline:
    """
    Complete automated retraining pipeline

    Orchestrates data collection, training, evaluation,
    and deployment of updated models
    """

    def __init__(self, config: RetrainingConfig):
        self.config = config
        self.logger = logging.getLogger(__name__)

        self.data_collector = DataCollector(
            data_path=Path(f"./data/{config.model_name}")
        )
        self.evaluator = PerformanceEvaluator()

        self.retraining_history: List[RetrainingEvent] = []
        self.last_retrain_time: Optional[datetime] = None

    def should_retrain(self,
                      current_metrics: Dict[str, float],
                      drift_score: float,
                      days_since_retrain: int) -> tuple[bool, RetrainingTrigger]:
        """Determine if retraining should be triggered"""

        # Check performance degradation
        baseline_metric = 0.85  # Load from model registry
        current_metric = current_metrics.get('f1', 0)

        degradation = baseline_metric - current_metric
        if degradation > self.config.performance_threshold:
            return True, RetrainingTrigger.PERFORMANCE_DEGRADATION

        # Check drift
        if drift_score > self.config.drift_threshold:
            return True, RetrainingTrigger.DATA_DRIFT

        # Check time-based schedule
        if days_since_retrain >= self.config.max_days_without_retrain:
            return True, RetrainingTrigger.SCHEDULED

        return False, None

    def execute_retraining(self,
                          trigger: RetrainingTrigger,
                          current_model: Any = None) -> RetrainingEvent:
        """Execute complete retraining workflow"""

        event_id = f"{self.config.model_name}_{datetime.now().isoformat()}"
        event = RetrainingEvent(
            event_id=event_id,
            timestamp=datetime.now(),
            trigger=trigger,
            model_version=f"v{len(self.retraining_history) + 1}",
            metrics={},
            data_stats={}
        )

        try:
            self.logger.info(
                f"Starting retraining for {self.config.model_name} "
                f"(trigger: {trigger.value})"
            )

            # Step 1: Collect training data
            datasets = self.data_collector.get_training_dataset(
                days_back=90,
                include_validation=True
            )

            event.data_stats = {
                'train_samples': len(datasets['train']),
                'validation_samples': len(datasets['validation'])
            }

            # Step 2: Train new model
            start_time = datetime.now()
            new_model = self._train_model(datasets['train'])
            event.training_duration = (datetime.now() - start_time).total_seconds()

            # Step 3: Evaluate new model
            new_metrics = self.evaluator.evaluate_model(
                new_model,
                datasets['validation']
            )
            event.metrics = new_metrics

            # Step 4: Compare with current model
            if current_model is not None:
                current_metrics = self.evaluator.evaluate_model(
                    current_model,
                    datasets['validation']
                )

                comparison = self.evaluator.compare_models(
                    new_metrics,
                    current_metrics
                )

                if comparison['recommendation'] != 'promote_new':
                    self.logger.warning(
                        "New model not better than current - skipping deployment"
                    )
                    event.success = False
                    event.error_message = "Model did not improve"
                    return event

            # Step 5: Save new model
            self._save_model(new_model, event.model_version)

            event.success = True
            self.last_retrain_time = datetime.now()

            self.logger.info(
                f"Retraining completed successfully: {event.model_version} "
                f"(F1: {new_metrics['f1']:.4f})"
            )

        except Exception as e:
            self.logger.error(f"Retraining failed: {str(e)}")
            event.success = False
            event.error_message = str(e)

        finally:
            self.retraining_history.append(event)

        return event

    def _train_model(self, train_data: pd.DataFrame) -> Any:
        """
        Train a new model using a simple, configurable baseline.

        Reference implementation: RandomForest with sensible defaults that can
        be overridden via RetrainingConfig.model_type if extended.
        """
        from sklearn.ensemble import RandomForestClassifier
        from sklearn.model_selection import train_test_split
        from sklearn.metrics import f1_score

        # Basic validation
        if "target" not in train_data.columns:
            raise ValueError("Training data must include 'target' column")

        X = train_data.drop(columns=["target"])
        y = train_data["target"]

        # Simple hold-out for quick sanity during training
        X_train, X_val, y_train, y_val = train_test_split(
            X, y, test_size=0.2, random_state=42, stratify=y
        )

        model = RandomForestClassifier(
            n_estimators=200,
            max_depth=None,
            min_samples_split=2,
            random_state=42,
            n_jobs=-1
        )
        model.fit(X_train, y_train)

        # Quick validation to fail fast on obviously bad models
        val_f1 = f1_score(y_val, model.predict(X_val), average="weighted")
        self.logger.info(f"Validation F1 during retraining: {val_f1:.4f}")

        return model

    def _save_model(self, model: Any, version: str):
        """Save model to storage"""

        model_path = Path(f"./models/{self.config.model_name}/{version}")
        model_path.mkdir(parents=True, exist_ok=True)

        with open(model_path / "model.pkl", 'wb') as f:
            pickle.dump(model, f)

        self.logger.info(f"Model saved: {model_path}")

# Example usage
config = RetrainingConfig(
    model_name="churn_predictor",
    model_type="classification",
    performance_threshold=0.05,
    drift_threshold=0.15,
    max_days_without_retrain=30
)

pipeline = AutomatedRetrainingPipeline(config)

# Check if retraining needed
current_metrics = {'f1': 0.78, 'accuracy': 0.82}
drift_score = 0.18
days_since_retrain = 15

should_retrain, trigger = pipeline.should_retrain(
    current_metrics,
    drift_score,
    days_since_retrain
)

if should_retrain:
    event = pipeline.execute_retraining(trigger)
    print(f"Retraining completed: {event.success}")
    print(f"New model metrics: {event.metrics}")
\end{lstlisting}

\section{Trigger Mechanisms for Retraining}

\subsection{Multi-Criteria Trigger System}

Production systems need sophisticated triggers that balance multiple signals to avoid both over-retraining (wasting resources) and under-retraining (allowing degradation).

\begin{lstlisting}[language=Python, caption=Advanced Retraining Trigger System]
from typing import Dict, List, Tuple, Optional
from dataclasses import dataclass
from datetime import datetime, timedelta
import numpy as np
from scipy import stats

@dataclass
class TriggerThreshold:
    """Threshold configuration for a trigger"""
    metric_name: str
    threshold_value: float
    comparison: str  # 'gt', 'lt', 'gte', 'lte'
    window_size: int = 1  # Number of consecutive violations
    severity: str = 'medium'  # 'low', 'medium', 'high', 'critical'

@dataclass
class TriggerEvaluation:
    """Result of trigger evaluation"""
    triggered: bool
    trigger_type: RetrainingTrigger
    severity: str
    confidence: float
    details: Dict[str, Any]
    timestamp: datetime = field(default_factory=datetime.now)

class RetrainingTriggerManager:
    """
    Manage complex retraining trigger logic

    Combines multiple signals (performance, drift, time, data volume)
    with configurable thresholds and cooldown periods
    """

    def __init__(self):
        self.logger = logging.getLogger(__name__)
        self.trigger_history: List[TriggerEvaluation] = []
        self.last_retrain_time: Optional[datetime] = None
        self.cooldown_hours = 24  # Minimum time between retrains

        # Configure trigger thresholds
        self.thresholds = {
            RetrainingTrigger.PERFORMANCE_DEGRADATION: [
                TriggerThreshold('accuracy', 0.05, 'gt', window_size=3, severity='high'),
                TriggerThreshold('f1_score', 0.07, 'gt', window_size=2, severity='critical'),
            ],
            RetrainingTrigger.DATA_DRIFT: [
                TriggerThreshold('psi_score', 0.2, 'gte', severity='medium'),
                TriggerThreshold('ks_statistic', 0.1, 'gte', severity='medium'),
            ],
            RetrainingTrigger.DATA_VOLUME: [
                TriggerThreshold('new_samples', 10000, 'gte', severity='low'),
            ],
        }

        # Metric history for windowed evaluation
        self.metric_history: Dict[str, List[float]] = {}

    def evaluate_triggers(self,
                         current_metrics: Dict[str, float],
                         baseline_metrics: Dict[str, float],
                         drift_metrics: Dict[str, float],
                         data_stats: Dict[str, int]) -> List[TriggerEvaluation]:
        """Evaluate all trigger conditions"""

        evaluations = []

        # Check if in cooldown period
        if self._in_cooldown():
            self.logger.info("In cooldown period - skipping trigger evaluation")
            return evaluations

        # 1. Performance degradation triggers
        perf_eval = self._evaluate_performance_triggers(
            current_metrics,
            baseline_metrics
        )
        if perf_eval:
            evaluations.append(perf_eval)

        # 2. Data drift triggers
        drift_eval = self._evaluate_drift_triggers(drift_metrics)
        if drift_eval:
            evaluations.append(drift_eval)

        # 3. Data volume triggers
        volume_eval = self._evaluate_volume_triggers(data_stats)
        if volume_eval:
            evaluations.append(volume_eval)

        # 4. Time-based triggers
        time_eval = self._evaluate_time_triggers()
        if time_eval:
            evaluations.append(time_eval)

        # Store trigger evaluations
        self.trigger_history.extend(evaluations)

        return evaluations

    def _in_cooldown(self) -> bool:
        """Check if we're in cooldown period after last retrain"""
        if self.last_retrain_time is None:
            return False

        time_since_retrain = datetime.now() - self.last_retrain_time
        return time_since_retrain < timedelta(hours=self.cooldown_hours)

    def _evaluate_performance_triggers(self,
                                      current: Dict[str, float],
                                      baseline: Dict[str, float]) -> Optional[TriggerEvaluation]:
        """Evaluate performance degradation triggers"""

        thresholds = self.thresholds[RetrainingTrigger.PERFORMANCE_DEGRADATION]
        violations = []

        for threshold in thresholds:
            metric = threshold.metric_name
            if metric not in current or metric not in baseline:
                continue

            degradation = baseline[metric] - current[metric]

            # Update metric history
            if metric not in self.metric_history:
                self.metric_history[metric] = []
            self.metric_history[metric].append(degradation)

            # Keep only window_size recent values
            self.metric_history[metric] = \
                self.metric_history[metric][-threshold.window_size:]

            # Check if threshold violated for entire window
            recent_values = self.metric_history[metric]
            if len(recent_values) >= threshold.window_size:
                if threshold.comparison == 'gt':
                    violated = all(v > threshold.threshold_value for v in recent_values)
                elif threshold.comparison == 'gte':
                    violated = all(v >= threshold.threshold_value for v in recent_values)
                elif threshold.comparison == 'lt':
                    violated = all(v < threshold.threshold_value for v in recent_values)
                else:  # lte
                    violated = all(v <= threshold.threshold_value for v in recent_values)

                if violated:
                    violations.append({
                        'metric': metric,
                        'degradation': np.mean(recent_values),
                        'threshold': threshold.threshold_value,
                        'severity': threshold.severity
                    })

        if violations:
            # Use highest severity
            max_severity = max(v['severity'] for v in violations)

            return TriggerEvaluation(
                triggered=True,
                trigger_type=RetrainingTrigger.PERFORMANCE_DEGRADATION,
                severity=max_severity,
                confidence=len(violations) / len(thresholds),
                details={'violations': violations}
            )

        return None

    def _evaluate_drift_triggers(self,
                                drift_metrics: Dict[str, float]) -> Optional[TriggerEvaluation]:
        """Evaluate data drift triggers"""

        thresholds = self.thresholds[RetrainingTrigger.DATA_DRIFT]
        violations = []

        for threshold in thresholds:
            metric = threshold.metric_name
            if metric not in drift_metrics:
                continue

            value = drift_metrics[metric]

            # Check threshold
            violated = False
            if threshold.comparison == 'gte':
                violated = value >= threshold.threshold_value
            elif threshold.comparison == 'gt':
                violated = value > threshold.threshold_value

            if violated:
                violations.append({
                    'metric': metric,
                    'value': value,
                    'threshold': threshold.threshold_value,
                    'severity': threshold.severity
                })

        if violations:
            max_severity = max(v['severity'] for v in violations)

            return TriggerEvaluation(
                triggered=True,
                trigger_type=RetrainingTrigger.DATA_DRIFT,
                severity=max_severity,
                confidence=min(1.0, len(violations) * 0.5),
                details={'violations': violations}
            )

        return None

    def _evaluate_volume_triggers(self,
                                 data_stats: Dict[str, int]) -> Optional[TriggerEvaluation]:
        """Evaluate data volume triggers"""

        thresholds = self.thresholds[RetrainingTrigger.DATA_VOLUME]

        for threshold in thresholds:
            if threshold.metric_name in data_stats:
                value = data_stats[threshold.metric_name]

                if threshold.comparison == 'gte' and value >= threshold.threshold_value:
                    return TriggerEvaluation(
                        triggered=True,
                        trigger_type=RetrainingTrigger.DATA_VOLUME,
                        severity=threshold.severity,
                        confidence=min(1.0, value / threshold.threshold_value),
                        details={'new_samples': value}
                    )

        return None

    def _evaluate_time_triggers(self) -> Optional[TriggerEvaluation]:
        """Evaluate time-based triggers"""

        if self.last_retrain_time is None:
            # Never retrained before
            return TriggerEvaluation(
                triggered=True,
                trigger_type=RetrainingTrigger.SCHEDULED,
                severity='low',
                confidence=1.0,
                details={'reason': 'initial_training'}
            )

        days_since_retrain = (datetime.now() - self.last_retrain_time).days

        # Monthly retraining schedule
        if days_since_retrain >= 30:
            return TriggerEvaluation(
                triggered=True,
                trigger_type=RetrainingTrigger.SCHEDULED,
                severity='medium',
                confidence=1.0,
                details={'days_since_retrain': days_since_retrain}
            )

        return None

    def should_retrain(self,
                      evaluations: List[TriggerEvaluation]) -> Tuple[bool, Optional[TriggerEvaluation]]:
        """Decide if retraining should proceed based on trigger evaluations"""

        if not evaluations:
            return False, None

        # Filter to triggered evaluations
        triggered = [e for e in evaluations if e.triggered]

        if not triggered:
            return False, None

        # Priority order: critical > high > medium > low
        severity_order = {'critical': 4, 'high': 3, 'medium': 2, 'low': 1}

        # Sort by severity and confidence
        triggered.sort(
            key=lambda x: (severity_order.get(x.severity, 0), x.confidence),
            reverse=True
        )

        primary_trigger = triggered[0]

        self.logger.info(
            f"Retraining triggered: {primary_trigger.trigger_type.value} "
            f"(severity: {primary_trigger.severity}, "
            f"confidence: {primary_trigger.confidence:.2f})"
        )

        return True, primary_trigger

    def record_retrain(self):
        """Record that retraining occurred"""
        self.last_retrain_time = datetime.now()

        # Clear metric history after retrain
        self.metric_history.clear()

# Example usage
trigger_manager = RetrainingTriggerManager()

# Simulate monitoring cycle
current_metrics = {
    'accuracy': 0.79,
    'f1_score': 0.76
}

baseline_metrics = {
    'accuracy': 0.85,
    'f1_score': 0.84
}

drift_metrics = {
    'psi_score': 0.25,
    'ks_statistic': 0.12
}

data_stats = {
    'new_samples': 12000
}

# Evaluate triggers
evaluations = trigger_manager.evaluate_triggers(
    current_metrics,
    baseline_metrics,
    drift_metrics,
    data_stats
)

# Check if should retrain
should_retrain, primary_trigger = trigger_manager.should_retrain(evaluations)

if should_retrain:
    print(f"Triggering retraining due to: {primary_trigger.trigger_type.value}")
    print(f"Severity: {primary_trigger.severity}")
    print(f"Details: {primary_trigger.details}")

    # Execute retraining...
    trigger_manager.record_retrain()
\end{lstlisting}

\section{Incremental Learning and Online Updating}

\subsection{Incremental Training Strategies}

Incremental learning allows models to learn from new data without full retraining, reducing computational costs and enabling faster adaptation.

\begin{lstlisting}[language=Python, caption=Incremental Learning System]
from typing import Dict, List, Optional, Any
import numpy as np
import pandas as pd
from sklearn.base import BaseEstimator
from datetime import datetime
import logging

class IncrementalLearner:
    """
    Manage incremental model updates

    Supports partial_fit for online learning algorithms and
    implements sliding window strategies for batch models
    """

    def __init__(self,
                 base_model: BaseEstimator,
                 update_strategy: str = 'sliding_window',
                 window_size: int = 10000,
                 update_frequency: int = 1000):
        """
        Initialize incremental learner

        Parameters:
        - base_model: Model supporting partial_fit or incremental training
        - update_strategy: 'sliding_window', 'cumulative', or 'temporal_weighting'
        - window_size: Max samples to keep in training buffer
        - update_frequency: Trigger update every N new samples
        """
        self.model = base_model
        self.update_strategy = update_strategy
        self.window_size = window_size
        self.update_frequency = update_frequency

        self.training_buffer: List[Dict[str, Any]] = []
        self.samples_since_update = 0
        self.total_updates = 0

        self.logger = logging.getLogger(__name__)

    def add_sample(self, features: Dict[str, Any], label: Any):
        """Add new labeled sample to training buffer"""

        self.training_buffer.append({
            'features': features,
            'label': label,
            'timestamp': datetime.now()
        })

        # Maintain window size
        if len(self.training_buffer) > self.window_size:
            if self.update_strategy == 'sliding_window':
                # Remove oldest sample
                self.training_buffer.pop(0)
            elif self.update_strategy == 'cumulative':
                # Keep all but downsample old data
                self._downsample_old_data()

        self.samples_since_update += 1

        # Check if update needed
        if self.samples_since_update >= self.update_frequency:
            self.update_model()

    def _downsample_old_data(self):
        """Downsample older data to maintain buffer size"""

        # Keep all recent data
        recent_cutoff = len(self.training_buffer) - self.window_size // 2
        recent_data = self.training_buffer[recent_cutoff:]

        # Sample from older data
        old_data = self.training_buffer[:recent_cutoff]
        old_sample_size = self.window_size - len(recent_data)

        if len(old_data) > old_sample_size:
            import random
            old_data = random.sample(old_data, old_sample_size)

        self.training_buffer = old_data + recent_data

    def update_model(self):
        """Update model with buffered samples"""

        if not self.training_buffer:
            return

        self.logger.info(
            f"Updating model with {len(self.training_buffer)} samples "
            f"(update #{self.total_updates + 1})"
        )

        # Prepare training data
        X = pd.DataFrame([s['features'] for s in self.training_buffer])
        y = pd.Series([s['label'] for s in self.training_buffer])

        # Apply temporal weighting if configured
        if self.update_strategy == 'temporal_weighting':
            sample_weights = self._calculate_temporal_weights()
        else:
            sample_weights = None

        # Update model
        if hasattr(self.model, 'partial_fit'):
            # Online learning algorithm
            if sample_weights is not None:
                self.model.partial_fit(X, y, sample_weight=sample_weights)
            else:
                self.model.partial_fit(X, y)
        else:
            # Retrain on windowed data
            if sample_weights is not None:
                self.model.fit(X, y, sample_weight=sample_weights)
            else:
                self.model.fit(X, y)

        self.total_updates += 1
        self.samples_since_update = 0

        self.logger.info(f"Model update completed (total updates: {self.total_updates})")

    def _calculate_temporal_weights(self) -> np.ndarray:
        """Calculate sample weights based on recency"""

        if not self.training_buffer:
            return np.array([])

        # Recent samples get higher weight
        ages = np.arange(len(self.training_buffer))
        weights = np.exp(-ages / (self.window_size / 3))

        # Normalize
        weights = weights / weights.sum() * len(weights)

        return weights

    def predict(self, features: Dict[str, Any]) -> Any:
        """Make prediction using current model"""
        X = pd.DataFrame([features])
        return self.model.predict(X)[0]

class OnlineModelUpdater:
    """
    Complete online model update system

    Handles continuous model updates with performance tracking
    and automatic rollback on degradation
    """

    def __init__(self,
                 model: BaseEstimator,
                 validation_window: int = 1000,
                 performance_threshold: float = 0.05):
        self.learner = IncrementalLearner(model, update_strategy='temporal_weighting')
        self.validation_window = validation_window
        self.performance_threshold = performance_threshold

        self.validation_buffer: List[Dict] = []
        self.performance_history: List[float] = []
        self.model_checkpoints: List[Any] = []

        self.logger = logging.getLogger(__name__)

    def process_prediction(self,
                          features: Dict[str, Any],
                          prediction: Any,
                          ground_truth: Optional[Any] = None):
        """Process prediction and optionally update model"""

        if ground_truth is not None:
            # Add to validation buffer
            self.validation_buffer.append({
                'features': features,
                'prediction': prediction,
                'ground_truth': ground_truth,
                'timestamp': datetime.now()
            })

            # Maintain validation window
            if len(self.validation_buffer) > self.validation_window:
                self.validation_buffer.pop(0)

            # Add to training buffer
            self.learner.add_sample(features, ground_truth)

            # Evaluate performance periodically
            if len(self.validation_buffer) >= self.validation_window:
                self._evaluate_performance()

    def _evaluate_performance(self):
        """Evaluate current model performance"""

        if len(self.validation_buffer) < 100:
            return

        # Calculate accuracy on validation buffer
        correct = sum(
            1 for sample in self.validation_buffer
            if sample['prediction'] == sample['ground_truth']
        )
        accuracy = correct / len(self.validation_buffer)

        self.performance_history.append(accuracy)

        # Check for degradation
        if len(self.performance_history) >= 2:
            recent_avg = np.mean(self.performance_history[-5:])
            baseline_avg = np.mean(self.performance_history[:10])

            degradation = baseline_avg - recent_avg

            if degradation > self.performance_threshold:
                self.logger.warning(
                    f"Performance degradation detected: {degradation:.4f} "
                    f"(recent: {recent_avg:.4f}, baseline: {baseline_avg:.4f})"
                )

                # Consider rollback or full retrain
                self._handle_degradation()

    def _handle_degradation(self):
        """Handle detected performance degradation"""

        self.logger.info("Initiating full retraining due to degradation")

        # In production, this would trigger automated retraining pipeline
        # For now, just log the event
        pass

# Example: Online learning for classification
from sklearn.naive_bayes import MultinomialNB
from sklearn.linear_model import SGDClassifier

# SGDClassifier supports partial_fit for online learning
online_model = SGDClassifier(
    loss='log_loss',
    learning_rate='adaptive',
    eta0=0.01,
    random_state=42
)

updater = OnlineModelUpdater(
    model=online_model,
    validation_window=1000,
    performance_threshold=0.05
)

# Simulate online predictions with delayed ground truth
for i in range(5000):
    features = {
        'feature_1': np.random.randn(),
        'feature_2': np.random.randn()
    }

    # Make prediction
    if i > 100:  # After initial training
        prediction = updater.learner.predict(features)
    else:
        prediction = 0

    # Simulate delayed ground truth (e.g., user feedback)
    ground_truth = np.random.choice([0, 1])

    # Process prediction with ground truth
    updater.process_prediction(features, prediction, ground_truth)

    if i % 1000 == 0:
        print(f"Processed {i} samples, total updates: {updater.learner.total_updates}")
\end{lstlisting}

\section{Concept Drift Handling Strategies}

\subsection{Detecting and Adapting to Concept Drift}

Concept drift occurs when the relationship between features and target changes over time, requiring model adaptation.

\begin{lstlisting}[language=Python, caption=Concept Drift Detection and Handling]
from typing import Dict, List, Tuple, Optional
import numpy as np
import pandas as pd
from scipy import stats
from collections import deque
from datetime import datetime, timedelta

class ConceptDriftDetector:
    """
    Detect concept drift in model predictions

    Uses multiple statistical tests to identify when the target
    distribution or feature-target relationship changes
    """

    def __init__(self,
                 reference_window: int = 1000,
                 detection_window: int = 500,
                 warning_threshold: float = 0.05,
                 drift_threshold: float = 0.01):
        self.reference_window = reference_window
        self.detection_window = detection_window
        self.warning_threshold = warning_threshold
        self.drift_threshold = drift_threshold

        self.reference_data: deque = deque(maxlen=reference_window)
        self.current_data: deque = deque(maxlen=detection_window)

        self.drift_detected = False
        self.warning_detected = False

        self.logger = logging.getLogger(__name__)

    def add_sample(self, features: np.ndarray, prediction: float, label: float):
        """Add new sample for drift detection"""

        sample = {
            'features': features,
            'prediction': prediction,
            'label': label,
            'error': abs(prediction - label),
            'timestamp': datetime.now()
        }

        # Add to current window
        self.current_data.append(sample)

        # Move to reference once current window is full
        if len(self.current_data) >= self.detection_window:
            if len(self.reference_data) < self.reference_window:
                # Still building reference
                self.reference_data.extend(list(self.current_data))
                self.current_data.clear()
            else:
                # Check for drift
                self._check_drift()

    def _check_drift(self) -> Tuple[bool, bool]:
        """Check for concept drift using multiple tests"""

        if len(self.reference_data) < self.reference_window:
            return False, False

        if len(self.current_data) < self.detection_window:
            return False, False

        # Extract errors from both windows
        ref_errors = np.array([s['error'] for s in self.reference_data])
        curr_errors = np.array([s['error'] for s in self.current_data])

        # 1. Kolmogorov-Smirnov test for error distribution
        ks_stat, ks_pvalue = stats.ks_2samp(ref_errors, curr_errors)

        # 2. Mean error comparison using t-test
        t_stat, t_pvalue = stats.ttest_ind(ref_errors, curr_errors)

        # 3. Page-Hinkley test for cumulative drift
        ph_drift = self._page_hinkley_test(curr_errors)

        # Evaluate results
        drift_signals = 0

        if ks_pvalue < self.drift_threshold:
            drift_signals += 1
            self.logger.warning(
                f"KS test detected distribution change: "
                f"statistic={ks_stat:.4f}, p-value={ks_pvalue:.6f}"
            )

        if t_pvalue < self.drift_threshold:
            drift_signals += 1
            self.logger.warning(
                f"T-test detected mean change: "
                f"ref_mean={ref_errors.mean():.4f}, "
                f"curr_mean={curr_errors.mean():.4f}, "
                f"p-value={t_pvalue:.6f}"
            )

        if ph_drift:
            drift_signals += 1
            self.logger.warning("Page-Hinkley test detected drift")

        # Determine drift status
        if drift_signals >= 2:
            self.drift_detected = True
            self.warning_detected = False
            self.logger.error(f"CONCEPT DRIFT DETECTED ({drift_signals} signals)")

            # Reset reference window to current data
            self._reset_reference_window()

            return True, False

        elif drift_signals == 1:
            self.warning_detected = True
            self.drift_detected = False
            self.logger.warning("Drift warning - monitoring closely")
            return False, True

        else:
            self.drift_detected = False
            self.warning_detected = False
            return False, False

    def _page_hinkley_test(self, errors: np.ndarray,
                          delta: float = 0.005,
                          lambda_: float = 50) -> bool:
        """
        Page-Hinkley test for detecting mean shifts

        Cumulative test that's sensitive to gradual drift
        """

        mean_error = errors.mean()
        cumsum = 0
        min_cumsum = 0

        for error in errors:
            cumsum += error - mean_error - delta

            if cumsum < min_cumsum:
                min_cumsum = cumsum

            ph_value = cumsum - min_cumsum

            if ph_value > lambda_:
                return True

        return False

    def _reset_reference_window(self):
        """Reset reference window after drift detected"""

        # Use current window as new reference
        self.reference_data.clear()
        self.reference_data.extend(list(self.current_data))
        self.current_data.clear()

        self.logger.info("Reference window reset with recent data")

class AdaptiveDriftHandler:
    """
    Handle concept drift with adaptive strategies

    Implements multiple adaptation strategies:
    - Ensemble of time-weighted models
    - Dynamic model selection
    - Gradual forgetting
    """

    def __init__(self, base_model_class, n_ensemble: int = 5):
        self.base_model_class = base_model_class
        self.n_ensemble = n_ensemble

        # Ensemble of models with different recency
        self.ensemble: List[Dict] = []
        self.drift_detector = ConceptDriftDetector()

        # Initialize ensemble
        for i in range(n_ensemble):
            self.ensemble.append({
                'model': base_model_class(),
                'age': 0,
                'weight': 1.0 / n_ensemble,
                'performance': []
            })

        self.logger = logging.getLogger(__name__)

    def train_ensemble(self, X: pd.DataFrame, y: pd.Series):
        """Train ensemble with temporal weighting"""

        n_samples = len(X)

        for i, member in enumerate(self.ensemble):
            # Each model trained on different time window
            window_size = n_samples // (i + 1)
            start_idx = max(0, n_samples - window_size)

            X_subset = X.iloc[start_idx:]
            y_subset = y.iloc[start_idx:]

            # Apply temporal weighting
            sample_weights = np.linspace(0.5, 1.0, len(X_subset))

            member['model'].fit(X_subset, y_subset, sample_weight=sample_weights)
            member['age'] = 0

            self.logger.info(
                f"Trained ensemble member {i} on {len(X_subset)} samples"
            )

    def predict(self, features: pd.DataFrame) -> np.ndarray:
        """Predict using weighted ensemble"""

        predictions = []

        for member in self.ensemble:
            pred = member['model'].predict(features)
            predictions.append(pred * member['weight'])

        # Weighted average
        final_pred = np.sum(predictions, axis=0)

        return final_pred

    def update_on_drift(self, X_new: pd.DataFrame, y_new: pd.Series):
        """Update ensemble when drift detected"""

        self.logger.info("Updating ensemble due to drift detection")

        # Retire oldest model
        self.ensemble.pop()

        # Add new model trained on recent data
        new_model = {
            'model': self.base_model_class(),
            'age': 0,
            'weight': 0.3,  # Higher weight for newest model
            'performance': []
        }
        new_model['model'].fit(X_new, y_new)

        self.ensemble.insert(0, new_model)

        # Reweight ensemble
        total_weight = sum(m['weight'] for m in self.ensemble)
        for member in self.ensemble:
            member['weight'] /= total_weight

        self.logger.info("Ensemble updated with new model")

    def adapt_weights(self, validation_errors: List[np.ndarray]):
        """Adapt ensemble weights based on recent performance"""

        if len(validation_errors) < len(self.ensemble):
            return

        # Calculate performance for each ensemble member
        for i, member in enumerate(self.ensemble):
            errors = validation_errors[i]
            performance = 1.0 / (1.0 + errors.mean())
            member['performance'].append(performance)

        # Update weights based on recent performance
        recent_performance = [
            np.mean(m['performance'][-10:]) if m['performance'] else 0.5
            for m in self.ensemble
        ]

        # Softmax weighting
        exp_perf = np.exp(np.array(recent_performance))
        new_weights = exp_perf / exp_perf.sum()

        for i, member in enumerate(self.ensemble):
            member['weight'] = new_weights[i]

        self.logger.info(f"Updated ensemble weights: {new_weights}")

# Example usage
from sklearn.ensemble import RandomForestClassifier

# Create adaptive drift handler
handler = AdaptiveDriftHandler(
    base_model_class=lambda: RandomForestClassifier(n_estimators=50),
    n_ensemble=3
)

# Simulate streaming data with concept drift
np.random.seed(42)

# Generate initial data
X_train = pd.DataFrame(np.random.randn(1000, 5))
y_train = (X_train.sum(axis=1) > 0).astype(int)

# Train ensemble
handler.train_ensemble(X_train, y_train)

# Simulate concept drift after 500 samples
for batch in range(10):
    # Generate batch with drift
    if batch < 5:
        # Original concept
        X_batch = pd.DataFrame(np.random.randn(100, 5))
        y_batch = (X_batch.sum(axis=1) > 0).astype(int)
    else:
        # Drifted concept (inverted relationship)
        X_batch = pd.DataFrame(np.random.randn(100, 5))
        y_batch = (X_batch.sum(axis=1) < 0).astype(int)

    # Make predictions
    predictions = handler.predict(X_batch)

    # Add samples to drift detector
    for i in range(len(X_batch)):
        handler.drift_detector.add_sample(
            X_batch.iloc[i].values,
            predictions[i],
            y_batch.iloc[i]
        )

    # Check for drift
    if handler.drift_detector.drift_detected:
        print(f"Drift detected at batch {batch}")
        handler.update_on_drift(X_batch, y_batch)
\end{lstlisting}

\subsection{Multi-Strategy Drift Adaptation}

Different types of drift require different adaptation strategies. A robust system should employ multiple approaches.

\begin{lstlisting}[language=Python, caption=Multi-Strategy Drift Adaptation Framework]
from typing import Dict, List, Optional, Any
from enum import Enum
import numpy as np
import pandas as pd
from datetime import datetime
import logging

class DriftType(Enum):
    """Types of drift that can occur"""
    SUDDEN = "sudden"  # Abrupt change
    GRADUAL = "gradual"  # Slow evolution
    INCREMENTAL = "incremental"  # Step-wise changes
    RECURRING = "recurring"  # Cyclical patterns
    NONE = "none"

class DriftAdaptationStrategy(Enum):
    """Adaptation strategies for different drift types"""
    RETRAIN_FULL = "retrain_full"
    RETRAIN_INCREMENTAL = "retrain_incremental"
    ENSEMBLE_UPDATE = "ensemble_update"
    WINDOW_SLIDE = "window_slide"
    NONE = "none"

class DriftCharacterizer:
    """
    Characterize the type and severity of drift

    Analyzes drift patterns to select appropriate adaptation strategy
    """

    def __init__(self, history_window: int = 100):
        self.history_window = history_window
        self.drift_history: List[Dict] = []
        self.logger = logging.getLogger(__name__)

    def characterize_drift(self,
                          error_history: List[float],
                          timestamps: List[datetime]) -> Dict[str, Any]:
        """Characterize drift type and severity"""

        if len(error_history) < 10:
            return {
                'drift_type': DriftType.NONE,
                'severity': 0.0,
                'recommended_strategy': DriftAdaptationStrategy.NONE
            }

        errors = np.array(error_history)

        # Calculate error rate of change
        error_diff = np.diff(errors)

        # Detect sudden drift (abrupt change)
        sudden_threshold = 3 * np.std(errors)
        sudden_changes = np.abs(error_diff) > sudden_threshold

        if np.any(sudden_changes):
            drift_type = DriftType.SUDDEN
            severity = np.max(np.abs(error_diff)) / np.mean(errors)
            strategy = DriftAdaptationStrategy.RETRAIN_FULL

        # Detect gradual drift (trend)
        elif self._detect_trend(errors):
            drift_type = DriftType.GRADUAL
            # Calculate trend severity
            from scipy.stats import linregress
            slope, _, _, _, _ = linregress(range(len(errors)), errors)
            severity = abs(slope) / np.mean(errors)
            strategy = DriftAdaptationStrategy.RETRAIN_INCREMENTAL

        # Detect recurring drift (seasonality)
        elif self._detect_seasonality(errors, timestamps):
            drift_type = DriftType.RECURRING
            severity = 0.5
            strategy = DriftAdaptationStrategy.ENSEMBLE_UPDATE

        else:
            drift_type = DriftType.NONE
            severity = 0.0
            strategy = DriftAdaptationStrategy.NONE

        result = {
            'drift_type': drift_type,
            'severity': severity,
            'recommended_strategy': strategy,
            'timestamp': datetime.now()
        }

        self.drift_history.append(result)

        self.logger.info(
            f"Drift characterization: {drift_type.value} "
            f"(severity: {severity:.3f}, strategy: {strategy.value})"
        )

        return result

    def _detect_trend(self, errors: np.ndarray,
                     significance: float = 0.05) -> bool:
        """Detect monotonic trend using Mann-Kendall test"""

        from scipy.stats import kendalltau

        n = len(errors)
        tau, p_value = kendalltau(range(n), errors)

        return p_value < significance and abs(tau) > 0.3

    def _detect_seasonality(self,
                           errors: np.ndarray,
                           timestamps: List[datetime]) -> bool:
        """Detect seasonal/recurring patterns"""

        if len(errors) < 50:
            return False

        # Simple autocorrelation check
        from scipy.signal import find_peaks

        # Compute autocorrelation
        autocorr = np.correlate(errors - errors.mean(),
                               errors - errors.mean(),
                               mode='full')
        autocorr = autocorr[len(autocorr)//2:]
        autocorr = autocorr / autocorr[0]

        # Find peaks in autocorrelation
        peaks, _ = find_peaks(autocorr[1:], height=0.3)

        # If significant peaks found, likely seasonal
        return len(peaks) > 0

class AdaptiveRetrainingOrchestrator:
    """
    Orchestrate adaptive retraining based on drift characteristics

    Selects and executes appropriate adaptation strategy
    """

    def __init__(self, base_model, config: RetrainingConfig):
        self.base_model = base_model
        self.config = config

        self.drift_characterizer = DriftCharacterizer()
        self.drift_detector = ConceptDriftDetector()
        self.adaptive_handler = AdaptiveDriftHandler(
            base_model_class=lambda: base_model.__class__()
        )

        self.error_history: List[float] = []
        self.timestamp_history: List[datetime] = []

        self.logger = logging.getLogger(__name__)

    def process_batch(self,
                     X: pd.DataFrame,
                     y_pred: np.ndarray,
                     y_true: np.ndarray):
        """Process batch of predictions and adapt if needed"""

        # Calculate batch error
        batch_error = np.mean(np.abs(y_pred - y_true))

        self.error_history.append(batch_error)
        self.timestamp_history.append(datetime.now())

        # Keep history bounded
        max_history = 1000
        if len(self.error_history) > max_history:
            self.error_history = self.error_history[-max_history:]
            self.timestamp_history = self.timestamp_history[-max_history:]

        # Check for drift
        for i in range(len(X)):
            self.drift_detector.add_sample(
                X.iloc[i].values,
                y_pred[i],
                y_true[i]
            )

        # If drift detected, characterize and adapt
        if self.drift_detector.drift_detected:
            self._handle_detected_drift(X, y_true)

    def _handle_detected_drift(self, X: pd.DataFrame, y: np.ndarray):
        """Handle detected drift with appropriate strategy"""

        # Characterize drift
        characterization = self.drift_characterizer.characterize_drift(
            self.error_history,
            self.timestamp_history
        )

        drift_type = characterization['drift_type']
        severity = characterization['severity']
        strategy = characterization['recommended_strategy']

        self.logger.warning(
            f"Executing adaptation strategy: {strategy.value} "
            f"for {drift_type.value} drift (severity: {severity:.3f})"
        )

        # Execute strategy
        if strategy == DriftAdaptationStrategy.RETRAIN_FULL:
            self._execute_full_retrain(X, y, severity)

        elif strategy == DriftAdaptationStrategy.RETRAIN_INCREMENTAL:
            self._execute_incremental_retrain(X, y)

        elif strategy == DriftAdaptationStrategy.ENSEMBLE_UPDATE:
            self._execute_ensemble_update(X, y)

        elif strategy == DriftAdaptationStrategy.WINDOW_SLIDE:
            self._execute_window_slide(X, y)

    def _execute_full_retrain(self, X: pd.DataFrame, y: np.ndarray,
                             severity: float):
        """Execute full model retraining"""

        self.logger.info("Executing full model retrain")

        # Collect larger dataset for full retrain
        # In production, fetch from data warehouse
        train_data = X  # Placeholder

        # Retrain from scratch
        self.base_model.fit(train_data, y)

        # Reset drift detector
        self.drift_detector = ConceptDriftDetector()

        self.logger.info("Full retrain completed")

    def _execute_incremental_retrain(self, X: pd.DataFrame, y: np.ndarray):
        """Execute incremental model update"""

        self.logger.info("Executing incremental retrain")

        if hasattr(self.base_model, 'partial_fit'):
            self.base_model.partial_fit(X, y)
        else:
            # For models without partial_fit, use warm start
            self.base_model.fit(X, y)

        self.logger.info("Incremental retrain completed")

    def _execute_ensemble_update(self, X: pd.DataFrame, y: np.ndarray):
        """Update ensemble for recurring drift"""

        self.logger.info("Executing ensemble update")

        self.adaptive_handler.update_on_drift(X, pd.Series(y))

        self.logger.info("Ensemble update completed")

    def _execute_window_slide(self, X: pd.DataFrame, y: np.ndarray):
        """Slide training window for continuous adaptation"""

        self.logger.info("Executing window slide")

        # Use only recent data
        window_size = self.config.min_data_points

        if len(X) > window_size:
            X_recent = X.iloc[-window_size:]
            y_recent = y[-window_size:]
        else:
            X_recent = X
            y_recent = y

        self.base_model.fit(X_recent, y_recent)

        self.logger.info("Window slide completed")

# Example: Comprehensive drift handling
from sklearn.ensemble import GradientBoostingRegressor

# Create orchestrator
model = GradientBoostingRegressor(n_estimators=100)
config = RetrainingConfig(
    model_name="adaptive_model",
    model_type="regression"
)

orchestrator = AdaptiveRetrainingOrchestrator(model, config)

# Simulate streaming data with different drift types
np.random.seed(42)

for batch_idx in range(20):
    # Generate data with evolving distribution
    if batch_idx < 5:
        # Normal data
        X = pd.DataFrame(np.random.randn(100, 5))
        y_true = X.sum(axis=1).values + np.random.randn(100) * 0.1

    elif batch_idx < 10:
        # Gradual drift
        shift = (batch_idx - 5) * 0.2
        X = pd.DataFrame(np.random.randn(100, 5) + shift)
        y_true = X.sum(axis=1).values + np.random.randn(100) * 0.1

    elif batch_idx < 15:
        # Sudden drift
        X = pd.DataFrame(np.random.randn(100, 5) * 2)
        y_true = X.sum(axis=1).values * -1 + np.random.randn(100) * 0.1

    else:
        # Back to normal (recurring)
        X = pd.DataFrame(np.random.randn(100, 5))
        y_true = X.sum(axis=1).values + np.random.randn(100) * 0.1

    # Make predictions
    if batch_idx > 0:
        y_pred = orchestrator.base_model.predict(X)
    else:
        # Initial training
        orchestrator.base_model.fit(X, y_true)
        y_pred = orchestrator.base_model.predict(X)

    # Process batch
    orchestrator.process_batch(X, y_pred, y_true)

    print(f"Batch {batch_idx}: Error = {np.mean(np.abs(y_pred - y_true)):.4f}")
\end{lstlisting}

\section{Model Freshness vs Stability Trade-offs}

\subsection{Balancing Update Frequency and Model Stability}

Frequent retraining keeps models fresh but can introduce instability, unexpected behavior changes, and increased costs. Finding the right balance requires understanding your application's requirements.

\begin{lstlisting}[language=Python, caption=Freshness-Stability Trade-off Manager]
from typing import Dict, List, Optional, Tuple
from dataclasses import dataclass
from datetime import datetime, timedelta
import numpy as np
import logging

@dataclass
class StabilityMetrics:
    """Track model stability over time"""
    prediction_variance: float  # Variance in predictions between versions
    behavioral_consistency: float  # How consistent model behavior is
    rollback_rate: float  # Rate of rollbacks due to issues
    deployment_success_rate: float  # Successful deployments
    avg_time_to_rollback: float  # Time to detect and rollback issues

@dataclass
class FreshnessMetrics:
    """Track model freshness indicators"""
    data_staleness_days: float  # Days since training data cutoff
    drift_score: float  # Current drift magnitude
    performance_gap: float  # Gap from baseline performance
    new_data_volume: int  # New samples since last retrain
    feature_coverage: float  # Coverage of feature space

class FreshnessStabilityManager:
    """
    Manage trade-off between model freshness and stability

    Adaptively adjusts retraining frequency based on
    application requirements and observed stability
    """

    def __init__(self,
                 min_retrain_interval_days: int = 7,
                 max_retrain_interval_days: int = 30,
                 stability_weight: float = 0.5):
        """
        Parameters:
        - min_retrain_interval_days: Minimum days between retrains
        - max_retrain_interval_days: Maximum days before forced retrain
        - stability_weight: Weight for stability vs freshness (0-1)
          0 = prioritize freshness, 1 = prioritize stability
        """
        self.min_interval = timedelta(days=min_retrain_interval_days)
        self.max_interval = timedelta(days=max_retrain_interval_days)
        self.stability_weight = stability_weight

        self.last_retrain_time: Optional[datetime] = None
        self.stability_history: List[StabilityMetrics] = []
        self.freshness_history: List[FreshnessMetrics] = []

        # Adaptive interval
        self.current_interval = timedelta(
            days=(min_retrain_interval_days + max_retrain_interval_days) / 2
        )

        self.logger = logging.getLogger(__name__)

    def should_retrain(self,
                      current_freshness: FreshnessMetrics,
                      current_stability: StabilityMetrics) -> Tuple[bool, str]:
        """Determine if retraining should occur"""

        if self.last_retrain_time is None:
            return True, "Initial training required"

        time_since_retrain = datetime.now() - self.last_retrain_time

        # Always retrain if max interval exceeded
        if time_since_retrain >= self.max_interval:
            return True, f"Maximum interval exceeded ({self.max_interval.days} days)"

        # Never retrain if within minimum interval
        if time_since_retrain < self.min_interval:
            return False, f"Within minimum interval ({self.min_interval.days} days)"

        # Calculate freshness and stability scores
        freshness_score = self._calculate_freshness_score(current_freshness)
        stability_score = self._calculate_stability_score(current_stability)

        # Weighted decision
        retrain_score = (
            (1 - self.stability_weight) * (1 - freshness_score) +  # Need for freshness
            self.stability_weight * stability_score  # Confidence in stability
        )

        # Adapt interval based on stability
        self._adapt_retraining_interval(current_stability)

        # Retrain if score indicates need
        threshold = 0.5
        should_retrain = retrain_score > threshold

        reason = (
            f"Retrain score: {retrain_score:.3f} "
            f"(freshness: {freshness_score:.3f}, stability: {stability_score:.3f})"
        )

        if should_retrain:
            self.logger.info(f"Retraining recommended: {reason}")
        else:
            self.logger.info(f"Skipping retrain: {reason}")

        return should_retrain, reason

    def _calculate_freshness_score(self, metrics: FreshnessMetrics) -> float:
        """Calculate freshness score (0=stale, 1=fresh)"""

        # Multiple freshness indicators
        scores = []

        # 1. Data staleness (0-1, lower is fresher)
        staleness_score = 1.0 - min(1.0, metrics.data_staleness_days / 30.0)
        scores.append(staleness_score)

        # 2. Drift score (0-1, lower is better)
        drift_score = 1.0 - min(1.0, metrics.drift_score / 0.5)
        scores.append(drift_score)

        # 3. Performance gap (0-1, smaller gap is better)
        perf_score = 1.0 - min(1.0, metrics.performance_gap / 0.1)
        scores.append(perf_score)

        # 4. New data availability (0-1, more data is better)
        data_score = min(1.0, metrics.new_data_volume / 10000.0)
        scores.append(data_score * 0.5)  # Lower weight

        # Average freshness
        return np.mean(scores)

    def _calculate_stability_score(self, metrics: StabilityMetrics) -> float:
        """Calculate stability score (0=unstable, 1=stable)"""

        scores = []

        # 1. Low prediction variance is good
        variance_score = 1.0 - min(1.0, metrics.prediction_variance / 0.1)
        scores.append(variance_score)

        # 2. High behavioral consistency is good
        scores.append(metrics.behavioral_consistency)

        # 3. Low rollback rate is good
        rollback_score = 1.0 - min(1.0, metrics.rollback_rate)
        scores.append(rollback_score)

        # 4. High deployment success is good
        scores.append(metrics.deployment_success_rate)

        return np.mean(scores)

    def _adapt_retraining_interval(self, stability: StabilityMetrics):
        """Adapt retraining interval based on stability"""

        # If system is stable, can retrain more frequently
        stability_score = self._calculate_stability_score(stability)

        if stability_score > 0.8:
            # High stability - can be more aggressive with updates
            self.current_interval = max(
                self.min_interval,
                self.current_interval - timedelta(days=1)
            )
        elif stability_score < 0.5:
            # Low stability - be more conservative
            self.current_interval = min(
                self.max_interval,
                self.current_interval + timedelta(days=2)
            )

        self.logger.info(
            f"Adapted retraining interval to {self.current_interval.days} days "
            f"based on stability score {stability_score:.3f}"
        )

    def record_retrain(self,
                      freshness: FreshnessMetrics,
                      stability: StabilityMetrics):
        """Record retraining event"""

        self.last_retrain_time = datetime.now()
        self.freshness_history.append(freshness)
        self.stability_history.append(stability)

        # Keep history bounded
        max_history = 100
        if len(self.stability_history) > max_history:
            self.stability_history = self.stability_history[-max_history:]
            self.freshness_history = self.freshness_history[-max_history:]

class PredictionStabilityTracker:
    """
    Track prediction stability across model versions

    Monitors how predictions change when models are updated
    """

    def __init__(self, sample_size: int = 1000):
        self.sample_size = sample_size
        self.reference_predictions: Optional[np.ndarray] = None
        self.reference_features: Optional[pd.DataFrame] = None
        self.logger = logging.getLogger(__name__)

    def set_reference(self, model, features: pd.DataFrame):
        """Set reference predictions from current model"""

        # Sample features for stability tracking
        if len(features) > self.sample_size:
            sample_indices = np.random.choice(
                len(features),
                self.sample_size,
                replace=False
            )
            self.reference_features = features.iloc[sample_indices].copy()
        else:
            self.reference_features = features.copy()

        # Get reference predictions
        self.reference_predictions = model.predict(self.reference_features)

        self.logger.info(
            f"Set reference predictions for {len(self.reference_features)} samples"
        )

    def measure_stability(self, new_model) -> Dict[str, float]:
        """Measure prediction stability of new model vs reference"""

        if self.reference_predictions is None:
            return {
                'prediction_variance': 0.0,
                'behavioral_consistency': 1.0,
                'avg_prediction_change': 0.0
            }

        # Get new predictions
        new_predictions = new_model.predict(self.reference_features)

        # Calculate stability metrics
        prediction_changes = np.abs(new_predictions - self.reference_predictions)

        # For classification, measure agreement
        if len(np.unique(self.reference_predictions)) < 10:
            agreement = np.mean(new_predictions == self.reference_predictions)
            behavioral_consistency = agreement
            avg_change = 1.0 - agreement
        else:
            # For regression, measure relative change
            relative_changes = prediction_changes / (
                np.abs(self.reference_predictions) + 1e-6
            )
            avg_change = np.mean(relative_changes)
            behavioral_consistency = 1.0 - min(1.0, avg_change)

        prediction_variance = np.var(prediction_changes)

        metrics = {
            'prediction_variance': float(prediction_variance),
            'behavioral_consistency': float(behavioral_consistency),
            'avg_prediction_change': float(np.mean(prediction_changes)),
            'max_prediction_change': float(np.max(prediction_changes)),
            'pct_changed_predictions': float(
                np.mean(prediction_changes > 0.01) * 100
            )
        }

        self.logger.info(
            f"Stability metrics: consistency={metrics['behavioral_consistency']:.3f}, "
            f"variance={metrics['prediction_variance']:.6f}"
        )

        return metrics

# Example: Managing freshness-stability trade-off
from sklearn.ensemble import RandomForestClassifier
import pandas as pd

# Create manager with balanced approach
manager = FreshnessStabilityManager(
    min_retrain_interval_days=7,
    max_retrain_interval_days=30,
    stability_weight=0.5  # Equal weight to freshness and stability
)

# Create stability tracker
stability_tracker = PredictionStabilityTracker()

# Initialize with reference model
reference_model = RandomForestClassifier(n_estimators=100, random_state=42)
X_ref = pd.DataFrame(np.random.randn(1000, 10))
y_ref = (X_ref.sum(axis=1) > 0).astype(int)
reference_model.fit(X_ref, y_ref)

stability_tracker.set_reference(reference_model, X_ref)

# Simulate model updates over time
manager.last_retrain_time = datetime.now() - timedelta(days=15)

for day in range(30):
    current_time = datetime.now() - timedelta(days=30-day)

    # Simulate freshness metrics
    freshness = FreshnessMetrics(
        data_staleness_days=(current_time - manager.last_retrain_time).days,
        drift_score=0.05 + (day * 0.01),  # Increasing drift
        performance_gap=0.02 + (day * 0.002),
        new_data_volume=day * 500,
        feature_coverage=0.95
    )

    # Train new model to check stability
    X_new = pd.DataFrame(np.random.randn(1000, 10) + day * 0.01)
    y_new = (X_new.sum(axis=1) > 0).astype(int)
    new_model = RandomForestClassifier(n_estimators=100, random_state=42+day)
    new_model.fit(X_new, y_new)

    # Measure stability
    stability_metrics_dict = stability_tracker.measure_stability(new_model)

    stability = StabilityMetrics(
        prediction_variance=stability_metrics_dict['prediction_variance'],
        behavioral_consistency=stability_metrics_dict['behavioral_consistency'],
        rollback_rate=0.05,
        deployment_success_rate=0.95,
        avg_time_to_rollback=2.0
    )

    # Check if should retrain
    should_retrain, reason = manager.should_retrain(freshness, stability)

    if should_retrain:
        print(f"Day {day}: RETRAINING - {reason}")
        manager.record_retrain(freshness, stability)
        # Update reference
        stability_tracker.set_reference(new_model, X_new)
    else:
        print(f"Day {day}: Skipping - {reason}")
\end{lstlisting}

\subsection{Cost-Aware Retraining Scheduling}

Retraining has real costs in compute resources, engineering time, and potential business disruption. A production system must optimize these costs.

\begin{lstlisting}[language=Python, caption=Cost-Aware Retraining Scheduler]
from typing import Dict, List, Optional
from dataclasses import dataclass
from datetime import datetime, timedelta, time
import numpy as np

@dataclass
class RetrainingCost:
    """Model retraining cost components"""
    compute_cost_usd: float  # Direct compute costs
    engineer_hours: float  # Human review time
    opportunity_cost_usd: float  # Lost revenue if model offline
    data_pipeline_cost_usd: float  # Cost to fetch/process data
    validation_cost_usd: float  # A/B test or canary costs

    @property
    def total_cost(self) -> float:
        return (self.compute_cost_usd +
                self.opportunity_cost_usd +
                self.data_pipeline_cost_usd +
                self.validation_cost_usd)

@dataclass
class RetrainingBenefit:
    """Expected benefits from retraining"""
    performance_improvement: float  # Expected metric improvement
    revenue_impact_usd: float  # Expected revenue impact
    risk_reduction: float  # Reduced risk from drift
    confidence: float  # Confidence in benefit estimates

class CostAwareScheduler:
    """
    Schedule retraining considering costs and benefits

    Optimizes timing and frequency to maximize ROI
    """

    def __init__(self,
                 cost_model: Dict[str, float],
                 business_hours: Tuple[time, time] = (time(9), time(17)),
                 preferred_days: List[int] = [1, 2, 3]):  # Tue, Wed, Thu
        """
        Parameters:
        - cost_model: Cost parameters for retraining operations
        - business_hours: Hours to avoid retraining (higher opportunity cost)
        - preferred_days: Preferred days of week (0=Monday, 6=Sunday)
        """
        self.cost_model = cost_model
        self.business_hours = business_hours
        self.preferred_days = preferred_days

        self.scheduled_retrains: List[datetime] = []
        self.logger = logging.getLogger(__name__)

    def estimate_retrain_cost(self,
                             training_samples: int,
                             estimated_duration_hours: float,
                             scheduled_time: datetime) -> RetrainingCost:
        """Estimate cost of retraining operation"""

        # Compute cost scales with data volume and duration
        base_compute_cost = self.cost_model.get('base_compute_cost', 10.0)
        cost_per_hour = self.cost_model.get('compute_cost_per_hour', 5.0)
        compute_cost = base_compute_cost + (cost_per_hour * estimated_duration_hours)

        # Data pipeline cost
        cost_per_sample = self.cost_model.get('cost_per_sample', 0.0001)
        data_cost = training_samples * cost_per_sample

        # Engineer oversight
        engineer_hourly_rate = self.cost_model.get('engineer_hourly_rate', 150.0)
        engineer_hours = 2.0  # Review and validation
        engineer_cost = engineer_hours * engineer_hourly_rate

        # Opportunity cost depends on timing
        opportunity_cost = self._calculate_opportunity_cost(
            scheduled_time,
            estimated_duration_hours
        )

        # Validation cost (canary deployment)
        validation_cost = self.cost_model.get('canary_cost', 20.0)

        return RetrainingCost(
            compute_cost_usd=compute_cost + engineer_cost,
            engineer_hours=engineer_hours,
            opportunity_cost_usd=opportunity_cost,
            data_pipeline_cost_usd=data_cost,
            validation_cost_usd=validation_cost
        )

    def _calculate_opportunity_cost(self,
                                   scheduled_time: datetime,
                                   duration_hours: float) -> float:
        """Calculate opportunity cost based on timing"""

        # Higher cost during business hours
        revenue_per_hour = self.cost_model.get('revenue_per_hour', 1000.0)
        downtime_fraction = self.cost_model.get('downtime_fraction', 0.1)

        # Check if deployment overlaps with business hours
        start_time = scheduled_time.time()
        business_start, business_end = self.business_hours

        if business_start <= start_time <= business_end:
            # Business hours - higher opportunity cost
            multiplier = 1.0
        else:
            # Off-hours - lower opportunity cost
            multiplier = 0.2

        opportunity_cost = (
            revenue_per_hour *
            duration_hours *
            downtime_fraction *
            multiplier
        )

        return opportunity_cost

    def estimate_retrain_benefit(self,
                                current_performance: float,
                                expected_improvement: float,
                                confidence: float = 0.7) -> RetrainingBenefit:
        """Estimate expected benefit from retraining"""

        # Convert performance improvement to revenue impact
        # This would be calibrated based on your business metrics
        revenue_per_point = self.cost_model.get('revenue_per_performance_point', 5000.0)

        revenue_impact = expected_improvement * revenue_per_point * confidence

        # Risk reduction value (avoiding drift-related failures)
        risk_reduction_value = self.cost_model.get('risk_reduction_value', 1000.0)

        return RetrainingBenefit(
            performance_improvement=expected_improvement,
            revenue_impact_usd=revenue_impact + risk_reduction_value,
            risk_reduction=0.2,  # 20% risk reduction
            confidence=confidence
        )

    def calculate_roi(self,
                     cost: RetrainingCost,
                     benefit: RetrainingBenefit,
                     time_horizon_days: int = 30) -> Dict[str, float]:
        """Calculate ROI for retraining decision"""

        # Project benefits over time horizon
        total_benefit = benefit.revenue_impact_usd

        # Subtract costs
        total_cost = cost.total_cost

        # Calculate ROI
        roi = ((total_benefit - total_cost) / total_cost) * 100 if total_cost > 0 else 0

        # Net present value (simplified)
        npv = total_benefit - total_cost

        return {
            'roi_percent': roi,
            'npv_usd': npv,
            'total_cost': total_cost,
            'total_benefit': total_benefit,
            'payback_days': (total_cost / total_benefit * time_horizon_days)
                           if total_benefit > 0 else float('inf')
        }

    def find_optimal_time(self,
                         earliest: datetime,
                         latest: datetime) -> datetime:
        """Find optimal time slot for retraining"""

        # Generate candidate times
        candidates = []
        current = earliest

        while current <= latest:
            # Prefer off-hours
            if current.time() < self.business_hours[0] or \
               current.time() > self.business_hours[1]:
                # Prefer specific days
                if current.weekday() in self.preferred_days:
                    score = 1.0
                else:
                    score = 0.5
            else:
                score = 0.1  # Discourage business hours

            candidates.append((current, score))
            current += timedelta(hours=1)

        # Sort by score and return best
        candidates.sort(key=lambda x: x[1], reverse=True)

        optimal_time = candidates[0][0] if candidates else earliest

        self.logger.info(
            f"Optimal retraining time: {optimal_time} "
            f"(score: {candidates[0][1]:.2f})"
        )

        return optimal_time

    def should_retrain_now(self,
                          current_performance: float,
                          expected_improvement: float,
                          training_samples: int,
                          duration_hours: float = 2.0) -> Tuple[bool, Dict]:
        """Decide if retraining should proceed based on ROI"""

        now = datetime.now()

        # Find optimal time
        optimal_time = self.find_optimal_time(
            earliest=now,
            latest=now + timedelta(days=7)
        )

        # Estimate costs and benefits
        cost = self.estimate_retrain_cost(
            training_samples,
            duration_hours,
            optimal_time
        )

        benefit = self.estimate_retrain_benefit(
            current_performance,
            expected_improvement
        )

        # Calculate ROI
        roi_metrics = self.calculate_roi(cost, benefit)

        # Decision: retrain if positive NPV and ROI > 100%
        should_retrain = (
            roi_metrics['npv_usd'] > 0 and
            roi_metrics['roi_percent'] > 100
        )

        decision = {
            'should_retrain': should_retrain,
            'optimal_time': optimal_time,
            'cost': cost,
            'benefit': benefit,
            'roi_metrics': roi_metrics,
            'recommendation': 'RETRAIN' if should_retrain else 'SKIP'
        }

        self.logger.info(
            f"Retraining decision: {decision['recommendation']} "
            f"(ROI: {roi_metrics['roi_percent']:.1f}%, NPV: ${roi_metrics['npv_usd']:.2f})"
        )

        return should_retrain, decision

# Example: Cost-aware scheduling
cost_model = {
    'base_compute_cost': 10.0,
    'compute_cost_per_hour': 5.0,
    'cost_per_sample': 0.0001,
    'engineer_hourly_rate': 150.0,
    'canary_cost': 20.0,
    'revenue_per_hour': 1000.0,
    'downtime_fraction': 0.1,
    'revenue_per_performance_point': 5000.0,
    'risk_reduction_value': 1000.0
}

scheduler = CostAwareScheduler(
    cost_model=cost_model,
    business_hours=(time(9), time(17)),
    preferred_days=[1, 2, 3]  # Tue-Thu
)

# Evaluate retraining decision
should_retrain, decision = scheduler.should_retrain_now(
    current_performance=0.85,
    expected_improvement=0.03,  # 3% improvement expected
    training_samples=50000,
    duration_hours=2.5
)

print(f"\nRetraining Decision: {decision['recommendation']}")
print(f"Optimal Time: {decision['optimal_time']}")
print(f"Total Cost: ${decision['cost'].total_cost:.2f}")
print(f"Expected Benefit: ${decision['benefit'].revenue_impact_usd:.2f}")
print(f"ROI: {decision['roi_metrics']['roi_percent']:.1f}%")
print(f"Payback Period: {decision['roi_metrics']['payback_days']:.1f} days")
\end{lstlisting}

\section{Continuous Learning Systems}

\subsection{Real-Time Model Adaptation}

Continuous learning systems maintain models that adapt in real-time to streaming data, enabling immediate response to changing patterns without batch retraining.

\begin{lstlisting}[language=Python, caption=Continuous Learning System Architecture]
from typing import Dict, List, Optional, Callable, Tuple
from dataclasses import dataclass, field
from datetime import datetime
from collections import deque
import numpy as np
import pandas as pd
from threading import Thread, Lock
from queue import Queue
import time
import logging

@dataclass
class LearningEvent:
    """Single learning event in continuous system"""
    features: Dict[str, Any]
    prediction: Optional[Any] = None
    ground_truth: Optional[Any] = None
    timestamp: datetime = field(default_factory=datetime.now)
    feedback_delay_seconds: Optional[float] = None

class ContinuousLearningBuffer:
    """
    Thread-safe buffer for continuous learning

    Manages streaming predictions and delayed ground truth
    """

    def __init__(self, max_buffer_size: int = 10000):
        self.max_buffer_size = max_buffer_size
        self.events: deque = deque(maxlen=max_buffer_size)
        self.lock = Lock()

        self.total_predictions = 0
        self.total_labels_received = 0

        self.logger = logging.getLogger(__name__)

    def add_prediction(self, event_id: str, features: Dict, prediction: Any):
        """Add prediction event"""

        with self.lock:
            event = LearningEvent(
                features=features,
                prediction=prediction,
                timestamp=datetime.now()
            )

            self.events.append((event_id, event))
            self.total_predictions += 1

    def add_ground_truth(self, event_id: str, ground_truth: Any) -> bool:
        """Add ground truth for existing prediction"""

        with self.lock:
            # Find matching prediction
            for i, (eid, event) in enumerate(self.events):
                if eid == event_id and event.ground_truth is None:
                    # Calculate feedback delay
                    delay = (datetime.now() - event.timestamp).total_seconds()
                    event.ground_truth = ground_truth
                    event.feedback_delay_seconds = delay

                    self.total_labels_received += 1

                    self.logger.debug(
                        f"Received ground truth for {event_id} "
                        f"(delay: {delay:.1f}s)"
                    )

                    return True

        return False

    def get_labeled_batch(self, batch_size: int = 100) -> List[Tuple[Dict, Any, Any]]:
        """Get batch of labeled examples for training"""

        labeled = []

        with self.lock:
            for event_id, event in self.events:
                if event.ground_truth is not None and event.prediction is not None:
                    labeled.append((
                        event.features,
                        event.prediction,
                        event.ground_truth
                    ))

                if len(labeled) >= batch_size:
                    break

        return labeled

    def get_stats(self) -> Dict[str, Any]:
        """Get buffer statistics"""

        with self.lock:
            labeled_count = sum(
                1 for _, event in self.events
                if event.ground_truth is not None
            )

            avg_delay = np.mean([
                event.feedback_delay_seconds
                for _, event in self.events
                if event.feedback_delay_seconds is not None
            ]) if labeled_count > 0 else 0

        return {
            'buffer_size': len(self.events),
            'total_predictions': self.total_predictions,
            'total_labels': self.total_labels_received,
            'labeled_in_buffer': labeled_count,
            'avg_feedback_delay_seconds': avg_delay
        }

class ContinuousLearningEngine:
    """
    Complete continuous learning engine

    Manages real-time model updates with performance monitoring
    and automatic quality control
    """

    def __init__(self,
                 model,
                 update_interval_seconds: int = 60,
                 min_batch_size: int = 100,
                 performance_check_interval: int = 1000):
        """
        Parameters:
        - model: Base model supporting partial_fit or incremental learning
        - update_interval_seconds: Time between model updates
        - min_batch_size: Minimum samples before update
        - performance_check_interval: Check performance every N predictions
        """
        self.model = model
        self.update_interval = update_interval_seconds
        self.min_batch_size = min_batch_size
        self.performance_check_interval = performance_check_interval

        self.buffer = ContinuousLearningBuffer()
        self.model_lock = Lock()

        # Performance tracking
        self.recent_errors: deque = deque(maxlen=1000)
        self.performance_history: List[Dict] = []

        # Control flags
        self.running = False
        self.update_thread: Optional[Thread] = None

        self.logger = logging.getLogger(__name__)

    def start(self):
        """Start continuous learning background thread"""

        if self.running:
            self.logger.warning("Continuous learning already running")
            return

        self.running = True
        self.update_thread = Thread(target=self._update_loop, daemon=True)
        self.update_thread.start()

        self.logger.info("Continuous learning started")

    def stop(self):
        """Stop continuous learning"""

        self.running = False

        if self.update_thread:
            self.update_thread.join(timeout=5)

        self.logger.info("Continuous learning stopped")

    def predict(self, event_id: str, features: Dict) -> Any:
        """Make prediction and store for later learning"""

        # Make prediction with current model
        with self.model_lock:
            X = pd.DataFrame([features])
            prediction = self.model.predict(X)[0]

        # Store in buffer
        self.buffer.add_prediction(event_id, features, prediction)

        # Check performance periodically
        if self.buffer.total_predictions % self.performance_check_interval == 0:
            self._check_performance()

        return prediction

    def provide_feedback(self, event_id: str, ground_truth: Any):
        """Provide ground truth feedback for prediction"""

        success = self.buffer.add_ground_truth(event_id, ground_truth)

        if not success:
            self.logger.warning(
                f"Could not find prediction for event {event_id}"
            )

    def _update_loop(self):
        """Background loop for continuous model updates"""

        self.logger.info("Update loop started")

        while self.running:
            try:
                # Wait for update interval
                time.sleep(self.update_interval)

                # Get labeled batch
                batch = self.buffer.get_labeled_batch(self.min_batch_size)

                if len(batch) >= self.min_batch_size:
                    self._update_model(batch)
                else:
                    self.logger.debug(
                        f"Insufficient labeled samples: {len(batch)} "
                        f"< {self.min_batch_size}"
                    )

            except Exception as e:
                self.logger.error(f"Error in update loop: {e}")

    def _update_model(self, batch: List[Tuple[Dict, Any, Any]]):
        """Update model with new batch"""

        self.logger.info(f"Updating model with {len(batch)} samples")

        # Prepare training data
        X = pd.DataFrame([features for features, _, _ in batch])
        y = pd.Series([label for _, _, label in batch])

        # Calculate current batch error
        predictions = [pred for _, pred, _ in batch]
        errors = [abs(pred - label) for pred, label in zip(predictions, y)]
        batch_error = np.mean(errors)

        self.recent_errors.extend(errors)

        # Update model
        try:
            with self.model_lock:
                if hasattr(self.model, 'partial_fit'):
                    self.model.partial_fit(X, y)
                else:
                    # Fall back to full fit (not ideal for continuous learning)
                    self.model.fit(X, y)

            self.logger.info(
                f"Model updated successfully (batch error: {batch_error:.4f})"
            )

        except Exception as e:
            self.logger.error(f"Model update failed: {e}")

    def _check_performance(self):
        """Check and log current performance"""

        if len(self.recent_errors) < 100:
            return

        stats = self.buffer.get_stats()

        avg_error = np.mean(list(self.recent_errors))
        error_std = np.std(list(self.recent_errors))

        performance = {
            'timestamp': datetime.now(),
            'avg_error': avg_error,
            'error_std': error_std,
            'predictions': stats['total_predictions'],
            'labels_received': stats['total_labels'],
            'label_rate': (stats['total_labels'] / stats['total_predictions']
                          if stats['total_predictions'] > 0 else 0)
        }

        self.performance_history.append(performance)

        self.logger.info(
            f"Performance check: avg_error={avg_error:.4f}, "
            f"label_rate={performance['label_rate']:.2%}"
        )

    def get_metrics(self) -> Dict[str, Any]:
        """Get current system metrics"""

        buffer_stats = self.buffer.get_stats()

        metrics = {
            'buffer': buffer_stats,
            'model_updates': len(self.performance_history),
            'recent_avg_error': (np.mean(list(self.recent_errors))
                                if self.recent_errors else 0),
            'running': self.running
        }

        return metrics

# Example: Continuous learning system
from sklearn.linear_model import SGDRegressor
import uuid

# Create continuous learning engine
model = SGDRegressor(
    learning_rate='adaptive',
    eta0=0.01,
    random_state=42
)

# Initial training
X_init = pd.DataFrame(np.random.randn(1000, 5))
y_init = X_init.sum(axis=1) + np.random.randn(1000) * 0.1
model.fit(X_init, y_init)

# Start continuous learning
engine = ContinuousLearningEngine(
    model=model,
    update_interval_seconds=5,  # Update every 5 seconds
    min_batch_size=50
)

engine.start()

# Simulate streaming predictions with delayed feedback
print("Starting continuous learning simulation...")

try:
    for i in range(500):
        # Generate features
        features = {f'feature_{j}': np.random.randn() for j in range(5)}

        # Make prediction
        event_id = str(uuid.uuid4())
        prediction = engine.predict(event_id, features)

        # Simulate delayed ground truth (e.g., user feedback)
        # In real system, this would come from user actions
        ground_truth = sum(features.values()) + np.random.randn() * 0.1

        # Simulate random delay in feedback (50% get feedback)
        if np.random.random() < 0.5:
            # Immediate feedback
            engine.provide_feedback(event_id, ground_truth)
        elif np.random.random() < 0.3:
            # Delayed feedback simulation (would be async in production)
            time.sleep(0.1)
            engine.provide_feedback(event_id, ground_truth)

        # Small delay between predictions
        time.sleep(0.01)

        # Print metrics periodically
        if i % 100 == 0:
            metrics = engine.get_metrics()
            print(f"\nIteration {i}:")
            print(f"  Buffer size: {metrics['buffer']['buffer_size']}")
            print(f"  Labels received: {metrics['buffer']['total_labels']}")
            print(f"  Model updates: {metrics['model_updates']}")
            print(f"  Recent error: {metrics['recent_avg_error']:.4f}")

finally:
    # Stop engine
    engine.stop()
    print("\nContinuous learning stopped")
\end{lstlisting}

\subsection{Champion-Challenger Framework}

Production continuous learning systems often use champion-challenger patterns to safely test new models while maintaining a stable baseline.

\begin{lstlisting}[language=Python, caption=Champion-Challenger Continuous Learning]
from typing import Dict, List, Optional, Any
from dataclasses import dataclass
from datetime import datetime, timedelta
import numpy as np
import pandas as pd
import logging

@dataclass
class ModelPerformance:
    """Track model performance metrics"""
    model_id: str
    predictions: int = 0
    correct: int = 0
    total_error: float = 0.0
    recent_errors: List[float] = field(default_factory=list)

    @property
    def accuracy(self) -> float:
        return self.correct / self.predictions if self.predictions > 0 else 0

    @property
    def avg_error(self) -> float:
        return self.total_error / self.predictions if self.predictions > 0 else 0

class ChampionChallengerSystem:
    """
    Manage multiple model versions with champion-challenger pattern

    Champion serves production traffic while challengers are evaluated
    on a small percentage to identify better models
    """

    def __init__(self,
                 champion_model,
                 challenger_traffic_pct: float = 0.1,
                 evaluation_window: int = 1000,
                 promotion_threshold: float = 0.02):
        """
        Parameters:
        - champion_model: Current production model
        - challenger_traffic_pct: Percentage of traffic for challengers
        - evaluation_window: Samples to evaluate before promotion decision
        - promotion_threshold: Min improvement required for promotion
        """
        self.champion = champion_model
        self.champion_id = "champion_v1"

        self.challengers: Dict[str, Any] = {}
        self.challenger_traffic_pct = challenger_traffic_pct
        self.evaluation_window = evaluation_window
        self.promotion_threshold = promotion_threshold

        # Performance tracking
        self.performance: Dict[str, ModelPerformance] = {
            self.champion_id: ModelPerformance(model_id=self.champion_id)
        }

        self.total_predictions = 0

        self.logger = logging.getLogger(__name__)

    def add_challenger(self, challenger_id: str, model):
        """Add new challenger model"""

        self.challengers[challenger_id] = model
        self.performance[challenger_id] = ModelPerformance(model_id=challenger_id)

        self.logger.info(f"Added challenger: {challenger_id}")

    def predict(self, features: Dict) -> Tuple[Any, str]:
        """
        Make prediction using champion or challenger

        Returns: (prediction, model_id)
        """

        self.total_predictions += 1

        # Decide which model to use
        if self.challengers and np.random.random() < self.challenger_traffic_pct:
            # Use random challenger
            challenger_id = np.random.choice(list(self.challengers.keys()))
            model = self.challengers[challenger_id]
            model_id = challenger_id
        else:
            # Use champion
            model = self.champion
            model_id = self.champion_id

        # Make prediction
        X = pd.DataFrame([features])
        prediction = model.predict(X)[0]

        return prediction, model_id

    def record_outcome(self, model_id: str, prediction: Any, ground_truth: Any):
        """Record prediction outcome for performance tracking"""

        if model_id not in self.performance:
            return

        perf = self.performance[model_id]
        perf.predictions += 1

        error = abs(prediction - ground_truth)
        perf.total_error += error
        perf.recent_errors.append(error)

        # Keep only recent errors
        if len(perf.recent_errors) > self.evaluation_window:
            perf.recent_errors = perf.recent_errors[-self.evaluation_window:]

        # For classification, track accuracy
        if prediction == ground_truth:
            perf.correct += 1

        # Check if should evaluate for promotion
        if perf.predictions % self.evaluation_window == 0:
            self._evaluate_promotion(model_id)

    def _evaluate_promotion(self, challenger_id: str):
        """Evaluate if challenger should be promoted to champion"""

        if challenger_id == self.champion_id:
            return

        challenger_perf = self.performance[challenger_id]
        champion_perf = self.performance[self.champion_id]

        # Need minimum samples for both
        if (challenger_perf.predictions < self.evaluation_window or
            champion_perf.predictions < self.evaluation_window):
            return

        # Compare performance
        challenger_error = np.mean(challenger_perf.recent_errors)
        champion_error = np.mean(champion_perf.recent_errors)

        improvement = champion_error - challenger_error

        self.logger.info(
            f"Evaluation: {challenger_id} vs {self.champion_id}: "
            f"challenger_error={challenger_error:.4f}, "
            f"champion_error={champion_error:.4f}, "
            f"improvement={improvement:.4f}"
        )

        # Promote if improvement exceeds threshold
        if improvement > self.promotion_threshold:
            self._promote_challenger(challenger_id)

    def _promote_challenger(self, challenger_id: str):
        """Promote challenger to champion"""

        self.logger.info(f"PROMOTING {challenger_id} to champion")

        # Old champion becomes archived
        old_champion_id = self.champion_id

        # Promote challenger
        self.champion = self.challengers[challenger_id]
        self.champion_id = challenger_id

        # Remove from challengers
        del self.challengers[challenger_id]

        # Reset performance tracking
        self.performance[challenger_id] = ModelPerformance(
            model_id=challenger_id
        )

    def get_leaderboard(self) -> List[Dict]:
        """Get current model performance leaderboard"""

        leaderboard = []

        for model_id, perf in self.performance.items():
            if perf.predictions == 0:
                continue

            status = "CHAMPION" if model_id == self.champion_id else "CHALLENGER"

            leaderboard.append({
                'model_id': model_id,
                'status': status,
                'predictions': perf.predictions,
                'avg_error': perf.avg_error,
                'recent_error': (np.mean(perf.recent_errors)
                               if perf.recent_errors else 0),
                'accuracy': perf.accuracy
            })

        # Sort by recent error
        leaderboard.sort(key=lambda x: x['recent_error'])

        return leaderboard

# Example: Champion-Challenger system
from sklearn.ensemble import GradientBoostingRegressor, RandomForestRegressor
from sklearn.linear_model import Ridge

# Create champion model
champion = GradientBoostingRegressor(n_estimators=50, random_state=42)
X_train = pd.DataFrame(np.random.randn(1000, 5))
y_train = X_train.sum(axis=1) + np.random.randn(1000) * 0.1
champion.fit(X_train, y_train)

# Create system
cc_system = ChampionChallengerSystem(
    champion_model=champion,
    challenger_traffic_pct=0.2,  # 20% to challengers
    evaluation_window=500,
    promotion_threshold=0.01
)

# Add challenger models
challenger1 = RandomForestRegressor(n_estimators=50, random_state=42)
challenger1.fit(X_train, y_train)
cc_system.add_challenger("random_forest_v1", challenger1)

challenger2 = Ridge(alpha=1.0)
challenger2.fit(X_train, y_train)
cc_system.add_challenger("ridge_v1", challenger2)

# Simulate streaming predictions
print("Testing Champion-Challenger system...\n")

for i in range(2000):
    # Generate sample
    features = {f'feature_{j}': np.random.randn() for j in range(5)}
    ground_truth = sum(features.values()) + np.random.randn() * 0.1

    # Make prediction
    prediction, model_id = cc_system.predict(features)

    # Record outcome
    cc_system.record_outcome(model_id, prediction, ground_truth)

    # Print leaderboard periodically
    if (i + 1) % 500 == 0:
        print(f"\n=== After {i+1} predictions ===")
        leaderboard = cc_system.get_leaderboard()

        for rank, entry in enumerate(leaderboard, 1):
            print(
                f"{rank}. {entry['status']:10s} {entry['model_id']:20s} "
                f"| Predictions: {entry['predictions']:4d} "
                f"| Recent Error: {entry['recent_error']:.4f}"
            )
\end{lstlisting}

\section{Practical Exercises}

\subsection{Exercise 1: Build an Automated Retraining Pipeline}

Implement a complete automated retraining pipeline with monitoring, triggers, and deployment.

\textbf{Requirements:}
\begin{itemize}
    \item Monitor model performance in production
    \item Detect when retraining is needed based on performance degradation
    \item Automatically collect training data from the past 90 days
    \item Train new model and evaluate against current champion
    \item Deploy new model if it shows >2\% improvement
    \item Send notifications on retraining events (success/failure)
    \item Log all retraining events with metrics and outcomes
\end{itemize}

\textbf{Dataset:} Use an e-commerce conversion prediction dataset with daily updates.

\textbf{Expected Output:}
\begin{itemize}
    \item Working retraining pipeline that runs on schedule or trigger
    \item Dashboard showing retraining history and model performance over time
    \item Alert system for retraining failures or performance issues
    \item Model registry tracking all versions with metadata
\end{itemize}

\textbf{Hints:}
\begin{itemize}
    \item Use the \texttt{AutomatedRetrainingPipeline} class as a starting point
    \item Integrate with a workflow orchestrator (Airflow, Prefect, or simple cron)
    \item Store models in versioned storage (S3, GCS, or local with git-lfs)
    \item Use email/Slack webhooks for notifications
    \item Track metrics in a database or file-based store
\end{itemize}

\subsection{Exercise 2: Implement Multi-Trigger Retraining Logic}

Create a sophisticated trigger system that combines multiple signals to make intelligent retraining decisions.

\textbf{Requirements:}
\begin{itemize}
    \item Implement performance degradation triggers (accuracy, precision, recall, F1)
    \item Add data drift triggers using PSI and KS statistics
    \item Include time-based triggers (monthly, or X days since last retrain)
    \item Implement data volume triggers (retrain when N new samples available)
    \item Add cooldown period to prevent excessive retraining
    \item Weight triggers by severity (critical, high, medium, low)
    \item Combine multiple trigger signals for final decision
\end{itemize}

\textbf{Expected Output:}
\begin{itemize}
    \item \texttt{RetrainingTriggerManager} class with configurable thresholds
    \item Trigger evaluation logs showing which conditions fired
    \item Visualization of trigger history and patterns
    \item Configuration file for easy threshold adjustment
\end{itemize}

\textbf{Test Scenarios:}
\begin{enumerate}
    \item Gradual performance degradation over 2 weeks
    \item Sudden concept drift (performance drops 10\% overnight)
    \item Accumulation of 50,000 new training samples
    \item 30 days elapsed since last retrain with stable performance
    \item Multiple triggers firing simultaneously
\end{enumerate}

\subsection{Exercise 3: Build an Incremental Learning System}

Implement an online learning system that updates continuously from streaming data.

\textbf{Requirements:}
\begin{itemize}
    \item Use streaming data with delayed feedback labels
    \item Implement sliding window strategy for training data
    \item Support both \texttt{partial\_fit} and full retraining approaches
    \item Track performance on hold-out validation set
    \item Implement temporal weighting (recent data weighted more)
    \item Add automatic rollback on performance degradation
    \item Monitor and log update frequency and batch sizes
\end{itemize}

\textbf{Models to Test:}
\begin{itemize}
    \item SGDClassifier (supports partial\_fit)
    \item SGDRegressor (supports partial\_fit)
    \item PassiveAggressiveClassifier
    \item MultinomialNB (for text classification)
\end{itemize}

\textbf{Expected Output:}
\begin{itemize}
    \item Incremental learning system processing streaming data
    \item Performance comparison: incremental vs. periodic full retraining
    \item Analysis of update frequency vs. model quality trade-off
    \item Memory usage profiling for different buffer sizes
\end{itemize}

\subsection{Exercise 4: Implement Concept Drift Detection}

Build a comprehensive drift detection system using multiple statistical tests.

\textbf{Requirements:}
\begin{itemize}
    \item Implement Kolmogorov-Smirnov test for distribution drift
    \item Add Page-Hinkley test for gradual drift
    \item Use ADWIN (Adaptive Windowing) for change detection
    \item Implement ensemble drift detector (combine multiple tests)
    \item Detect both feature drift and concept drift
    \item Characterize drift type (sudden, gradual, recurring)
    \item Trigger appropriate adaptation strategy based on drift type
\end{itemize}

\textbf{Test Datasets:}
\begin{enumerate}
    \item Synthetic data with known drift points
    \item Weather prediction with seasonal patterns
    \item Financial time series with market regime changes
    \item User behavior data with gradual trend shifts
\end{enumerate}

\textbf{Expected Output:}
\begin{itemize}
    \item Drift detection accuracy (precision/recall for drift events)
    \item Time-to-detection for different drift types
    \item False positive rate analysis
    \item Comparison of different drift detection methods
    \item Visualization showing detected drift points
\end{itemize}

\subsection{Exercise 5: Optimize Freshness-Stability Trade-off}

Develop a system that adaptively balances model freshness against stability requirements.

\textbf{Requirements:}
\begin{itemize}
    \item Track both freshness metrics (data age, drift score) and stability metrics (prediction variance, rollback rate)
    \item Implement adaptive retraining interval that adjusts based on stability
    \item Add prediction stability tracker that compares model versions
    \item Create configurable freshness-stability weight parameter
    \item Log all retraining decisions with justifications
    \item Measure business impact of different configurations
\end{itemize}

\textbf{Scenarios to Test:}
\begin{enumerate}
    \item High-churn environment (prioritize freshness)
    \item Regulated environment (prioritize stability)
    \item Balanced scenario (equal weight)
    \item Varying data drift rates over time
\end{enumerate}

\textbf{Metrics to Track:}
\begin{itemize}
    \item Average model age (staleness)
    \item Prediction variance between model versions
    \item Number of rollbacks required
    \item Overall model performance over time
    \item Retraining frequency and costs
\end{itemize}

\textbf{Expected Output:}
\begin{itemize}
    \item Optimal freshness-stability configuration for each scenario
    \item Analysis showing trade-off curves
    \item Recommendation system for weight selection
    \item Cost-benefit analysis of different strategies
\end{itemize}

\subsection{Exercise 6: Build a Cost-Aware Retraining Scheduler}

Create an intelligent scheduler that optimizes retraining timing for cost and business impact.

\textbf{Requirements:}
\begin{itemize}
    \item Estimate retraining costs (compute, data pipeline, engineering)
    \item Calculate expected benefits (performance improvement, revenue impact)
    \item Compute ROI for retraining decisions
    \item Find optimal time windows (avoid business hours, prefer off-peak)
    \item Consider opportunity costs of model downtime
    \item Implement multi-day lookahead scheduling
    \item Add budget constraints and cost tracking
\end{itemize}

\textbf{Cost Model Components:}
\begin{itemize}
    \item Compute cost: based on training time and instance type
    \item Data pipeline cost: based on data volume processed
    \item Engineering oversight: human review and validation time
    \item Opportunity cost: potential revenue impact during deployment
    \item Validation cost: A/B testing or canary deployment costs
\end{itemize}

\textbf{Expected Output:}
\begin{itemize}
    \item Scheduler that recommends optimal retraining times
    \item ROI calculations for each retraining decision
    \item Cost tracking dashboard
    \item Schedule optimization showing cost savings vs. naive scheduling
    \item Budget forecasting for monthly retraining costs
\end{itemize}

\subsection{Exercise 7: Implement a Continuous Learning System with Champion-Challenger}

Build a production-ready continuous learning system with safe model updates.

\textbf{Requirements:}
\begin{itemize}
    \item Implement thread-safe continuous learning buffer
    \item Handle delayed feedback labels (predictions now, labels later)
    \item Run background model update thread
    \item Implement champion-challenger pattern for safe updates
    \item Route small percentage of traffic to challenger models
    \item Automatically promote challenger if performance improves
    \item Include automatic rollback on degradation
    \item Track and compare performance of all model versions
\end{itemize}

\textbf{System Components:}
\begin{enumerate}
    \item \textbf{Prediction Service:} Serves real-time predictions, tracks events
    \item \textbf{Feedback Collector:} Receives delayed ground truth labels
    \item \textbf{Learning Buffer:} Thread-safe storage for training events
    \item \textbf{Update Engine:} Periodic model updates from buffered data
    \item \textbf{Model Registry:} Manages champion and challenger models
    \item \textbf{Traffic Router:} Splits traffic between models for evaluation
    \item \textbf{Performance Monitor:} Tracks and compares model performance
    \item \textbf{Promotion Logic:} Decides when to promote challengers
\end{enumerate}

\textbf{Expected Output:}
\begin{itemize}
    \item Working continuous learning system processing streaming data
    \item Demonstration of automatic model promotion
    \item Performance comparison: champion vs. challengers over time
    \item Analysis of different traffic split percentages
    \item Latency and throughput benchmarks
    \item Safety analysis: rollback time, degradation detection delay
\end{itemize}

\textbf{Advanced Extensions:}
\begin{itemize}
    \item Multi-armed bandit for dynamic traffic allocation
    \item Bayesian optimization for challenger hyperparameters
    \item Distributed learning across multiple workers
    \item Integration with feature store for consistent features
    \item A/B test statistical significance testing
\end{itemize}

\section{Summary}

Model maintenance and retraining are critical for keeping ML systems healthy in production. This chapter covered:

\textbf{Automated Retraining Pipelines:}
\begin{itemize}
    \item Complete pipeline architecture with data collection, training, evaluation, and deployment
    \item Automated orchestration and error handling
    \item Integration with model registries and deployment systems
\end{itemize}

\textbf{Intelligent Trigger Mechanisms:}
\begin{itemize}
    \item Multi-criteria triggers combining performance, drift, time, and data volume signals
    \item Configurable thresholds with severity levels
    \item Cooldown periods to prevent excessive retraining
    \item Trigger evaluation and decision logic
\end{itemize}

\textbf{Incremental and Online Learning:}
\begin{itemize}
    \item Streaming data processing with delayed feedback
    \item Sliding window and temporal weighting strategies
    \item Online learning algorithms using partial\_fit
    \item Performance monitoring and automatic rollback
\end{itemize}

\textbf{Concept Drift Handling:}
\begin{itemize}
    \item Statistical drift detection methods (KS test, Page-Hinkley, ADWIN)
    \item Drift characterization (sudden, gradual, recurring)
    \item Adaptive strategies for different drift types
    \item Ensemble approaches for robust adaptation
\end{itemize}

\textbf{Freshness vs. Stability Trade-offs:}
\begin{itemize}
    \item Balancing model currency against behavioral consistency
    \item Adaptive retraining intervals based on stability
    \item Prediction stability tracking across versions
    \item Application-specific optimization
\end{itemize}

\textbf{Cost-Aware Scheduling:}
\begin{itemize}
    \item ROI-based retraining decisions
    \item Cost modeling for compute, data, and opportunity costs
    \item Optimal timing to minimize business impact
    \item Budget constraints and forecasting
\end{itemize}

\textbf{Continuous Learning Systems:}
\begin{itemize}
    \item Real-time model adaptation from streaming data
    \item Thread-safe buffering and background updates
    \item Champion-challenger pattern for safe deployment
    \item Automatic promotion based on performance
\end{itemize}

\textbf{Key Takeaways:}

\begin{enumerate}
    \item \textbf{Automate Everything:} Manual retraining doesn't scale. Build automated pipelines with proper monitoring and error handling.

    \item \textbf{Use Multiple Signals:} Don't rely on a single trigger. Combine performance metrics, drift detection, data volume, and time-based schedules.

    \item \textbf{Start Conservative:} Begin with longer retraining intervals and stricter promotion criteria. Relax constraints as you gain confidence in your system.

    \item \textbf{Monitor Continuously:} Track both model performance and system health. Set up alerts for anomalies and failures.

    \item \textbf{Handle Drift Adaptively:} Different types of drift require different responses. Characterize drift patterns and adapt accordingly.

    \item \textbf{Balance Trade-offs:} Consider freshness vs. stability, cost vs. performance, and automation vs. control. Optimize for your specific business needs.

    \item \textbf{Test Safely:} Use champion-challenger patterns, canary deployments, and automatic rollbacks to minimize risk.

    \item \textbf{Measure Business Impact:} Connect retraining decisions to business metrics. Calculate ROI to justify infrastructure investment.
\end{enumerate}

Production ML systems require ongoing maintenance. The patterns and code in this chapter provide a foundation for building robust, automated retraining systems that keep your models performing well as the world changes.

